\begin{tikzpicture}[scale=1.1, >=Stealth]
    % --- Riquadro Sinistro: Circonferenza Unitaria ---
    \begin{scope}[shift={(0,0)}]
        % Titolo
        \node[font=\bfseries\color{navyblue}] at (0, 2.6) {Circonferenza: $X^2 + Y^2 = 1$};
        
        % Assi
        \draw[->, thick, gray!80] (-2.2, 0) -- (2.2, 0) node[right] {$X$};
        \draw[->, thick, gray!80] (0, -2.2) -- (0, 2.2) node[above] {$Y$};
        
        % Settore circolare ombreggiato (Area = t/2)
        \def\angcircle{48}
        \fill[coolblue!25] (0,0) -- (1.8,0) arc (0:\angcircle:1.8) -- cycle;
        \node[coolblue!90!black, font=\bfseries] at (24:1.0) {Area $= \frac{t}{2}$};

        % Circonferenza
        \draw[very thick, coolblue] (0,0) circle (1.8);

        % Raggio e coordinate
        \coordinate (Oc) at (0,0);
        \coordinate (Pc) at ({\angcircle}:1.8);
        \draw[thick, crimsonred] (Oc) -- (Pc);
        \fill[crimsonred] (Pc) circle (2pt);
        \node[above right, crimsonred, font=\footnotesize] at (Pc) {$P(\cos t, \sin t)$};

        % Proiezioni
        \draw[dashed, gray] (Pc) -- ({\angcircle}:1.8 |- Oc) node[below, font=\scriptsize] {$\cos t$};
        \draw[dashed, gray] (Pc) -- (Pc -| Oc) node[left, font=\scriptsize] {$\sin t$};

        % Angolo t
        \draw[->, thick, forestgreen] (0.6,0) arc (0:\angcircle:0.6);
        \node[forestgreen, font=\footnotesize] at (24:0.8) {$t$};
    \end{scope}

    % --- Riquadro Destro: Iperbole Equilatera ---
    \begin{scope}[shift={(6.5,0)}]
        % Titolo
        \node[font=\bfseries\color{navyblue}] at (0, 2.6) {Iperbole Equilatera: $X^2 - Y^2 = 1$};
        
        % Assi
        \draw[->, thick, gray!80] (-1.2, 0) -- (3.2, 0) node[right] {$X$};
        \draw[->, thick, gray!80] (0, -2.2) -- (0, 2.2) node[above] {$Y$};

        % Asintoti Y = +- X
        \draw[dashed, gray!70] (-1.2, -1.2) -- (2.2, 2.2) node[above right, font=\tiny] {$Y=X$};
        \draw[dashed, gray!70] (-1.2, 1.2) -- (2.2, -2.2) node[below right, font=\tiny] {$Y=-X$};

        % Settore iperbolico ombreggiato (Area = t/2)
        % Usiamo coordinate parametriche per riempire: t_val = 1.1, scala = 1.2
        % x = a cosh(u), y = a sinh(u) con a = 1.4 (scala grafica)
        \def\a{1.3}
        \def\tmax{0.95}
        \fill[forestgreen!25] (0,0) -- ({\a*cosh(0)},0)
            \foreach \u in {0.05, 0.1, ..., \tmax} {
                -- ({\a*cosh(\u)}, {\a*sinh(\u)})
            }
            -- (0,0);
        \node[forestgreen!90!black, font=\bfseries] at (1.1, 0.4) {Area $= \frac{t}{2}$};

        % Ramo destro iperbole
        \draw[very thick, forestgreen, domain=-1.2:1.2, samples=100] 
            plot ({\a*cosh(\x)}, {\a*sinh(\x)});

        % Punto P e raggio
        \coordinate (Oh) at (0,0);
        \coordinate (Ph) at ({\a*cosh(\tmax)}, {\a*sinh(\tmax)});
        \draw[thick, crimsonred] (Oh) -- (Ph);
        \fill[crimsonred] (Ph) circle (2pt);
        \node[above right, crimsonred, font=\footnotesize] at (Ph) {$P(\cosh t, \sinh t)$};

        % Proiezioni
        \draw[dashed, gray] (Ph) -- ({\a*cosh(\tmax)}, 0) node[below, font=\scriptsize] {$\cosh t$};
        \draw[dashed, gray] (Ph) -- (0, {\a*sinh(\tmax)}) node[left, font=\scriptsize] {$\sinh t$};

        % Vertice
        \fill[black] (\a, 0) circle (1.5pt) node[below left, font=\scriptsize] {$(1,0)$};
    \end{scope}
\end{tikzpicture}
